\section{Other tests and the remaining questions}\label{sec:conclusions}
\subsection{A recurrence can preserve too much}
The finite complete bases in Section~\ref{sec:linear} suggested
transporting them between neighboring indices. There is an elementary
test that such a proposal must pass before estimating transition
matrices. Suppose fixed rational rows $v_1,\ldots,v_r$ span
$(1,0,\ldots,0)$, and every nonzero integer row $c_{i,n}$ is proportional
to its corresponding $v_i$. Replace each $v_i$ by its primitive integer
representative $u_i$. Then $c_{i,n}=k_{i,n}u_i$ with nonzero integral
$k_{i,n}$. For $\theta=(1,\theta_1,\ldots)$, at least one fixed
$u_i\theta$ is nonzero, and
\[
 |c_{i,n}\theta|\geq|u_i\theta|>0.
\]
Such transport cannot make all the values tend to zero. The proposed
neighboring-index update preserved precisely these individual
coefficient directions. A recurrence that supplies new combinations
or an additional quantitative constraint is not excluded; merely
propagating the same rational lines cannot solve the height problem.

\subsection{A Loewner specialization with an inconclusive cost comparison}
We also considered the divided-difference construction of Litsey
\cite{Litsey}, specialized to
\[
 F(z)=\frac{\psi''(z+1/4)-\psi''(z+3/4)}{1536},
 \qquad F'(0)=\bet.
\]
Here $\psi=\Gamma'/\Gamma$ is the digamma function.
The two cubic pole classes of $F$ give a degree-six difference equation.
The specialization of the cited nonvanishing argument applies for
$r\geq4$, $m\geq0$ and every rational evaluation point: some minor
deleting at most three common nodes is nonzero. The opposite signs
of the two pole classes do not invalidate that argument.

One fixed geometry, $r=200k$, $m=13k$, was assessed on the vertical
contour $\Re z=-1/8$, between the nonnegative nodes and the negative
poles. It avoids enclosed-pole terms. The
moving-contour and deletion errors are uniform lower-order terms,
including every position of the at most three deletions. On the
$r^2$ scale the certified lower decay rate was greater than $2.488$,
while the selected common clearer had upper rate $19.2552$.
These bounds establish no useful comparison. Neither is asserted to
be optimal or necessary; the outcome is inconclusive. This assessment
is retained in the record catalog rather than promoted to a failure
theorem for the Loewner construction.

\subsection{What would change the conclusions}
The results leave three different kinds of question.
For the linear forms, uniform formal rank is established and eleven
directions have sufficient height control; the missing condition is a
complete useful collection, or nonvanishing for the same small forms.
For the thirteen-value Gram construction, the residual analytic scale
is known, and the missing input is its primitive arithmetic.
For the low-degree mixed-pole determinants, genuine uniform content
savings are known, but the true arithmetic and achieved decay have not
been compared sharply enough.

The final substitution result is different. Its necessary dyadic cost
and the uniform lower bound for its error envelope prove that the
specified certificate fails. Improving a convenient entry clearer
cannot evade that theorem. A successful modification would have to
change an assumption of the certificate, such as use of the global
polynomial clearer or the generic $q^{-2}$ error envelope, and then
establish the corresponding new arithmetic and approximation estimates.

These distinctions explain the value of retaining both successful
lemmas and unsuccessful constructions. They identify which parts can
be reused and which estimates would actually move the argument forward.

\subsection{Prepared with Codex}
Prepared with OpenAI Codex, using GPT-6-Astra at maximum reasoning effort.
The cooperating sessions developed and cross-checked the underlying
records; this manuscript assembles their mathematical results.

\subsection{Verification}
The manuscript passed a separate cooperating-session mathematical and
exposition review. The accompanying records state the review and finite
checks' scope; neither constitutes formal verification.
